
\section*{Глава 3. Разработка веб-приложения}
\addcontentsline{toc}{section}{Глава 3. Разработка веб-приложения}
\stepcounter{section}

\subsection{Разработка базовой структуры приложения и вёрстка}

В соответствии с задачей 3.1, была реализована компонентная структура фронтенд-части
на базе библиотеки \textbf{React.js}. Использование хуков (\textit{Hooks}) позволило
эффективно управлять состоянием формы заказа. Пример базового компонента представлен
в листинге~\ref{lst:component}.

\begin{lstlisting}[caption={Пример структуры компонента карточки заказа}, label={lst:component}]
// Комментарий 
const OrderCard = ({ order }) => {
  return (
    <div className="order-card">
      <h3>Заказ #{order.id}</h3>
      <p>Адрес: {order.address}</p>
      <button onClick={() => addToRoute(order.id)}>
        Добавить в маршрут
      </button>
    </div>
  );
};
\end{lstlisting}

\subsection{Реализация модуля авторизации}

Согласно задаче 3.2, в системе реализована ролевая модель доступа. Для обеспечения
безопасности используется протокол \textbf{JWT} (JSON Web Token). При входе в
систему сервер генерирует токен, содержащий поле \texttt{role}, которое может
принимать значения \texttt{'dispatcher'} или \texttt{'courier'}.

\subsection{Интеграция картографических сервисов}

Ключевым этапом разработки (задача 3.4) стала интеграция API \linebreak \textbf{OpenStreetMap}
через библиотеку \textit{Leaflet}. Для решения транспортной задачи был выбран
жадный алгоритм с последующей оптимизацией методом локального поиска.

Данные о маршруте передаются в формате JSON (см. листинг~\ref{lst:route_json}).

\begin{lstlisting}[float, floatplacement=ht!, caption={Структура данных оптимизированного маршрута}, label={lst:route_json}]
{
  "route_id": "r-101",
  "courier_id": "u-45",
  "waypoints": [
    {"lat": 55.75, "lng": 37.61, "sequence": 1},
    {"lat": 55.80, "lng": 37.70, "sequence": 2}
  ],
  "total_distance_km": 12.4
}
\end{lstlisting}

Как видно из листинга~\ref{lst:route_json}, каждая точка содержит порядковый
номер (\texttt{sequence}), что позволяет курьеру следовать по оптимальному пути,
рассчитанному на сервере.
